\section{A typed calculus for the forbidden configuration}
\label{sec:enc-wrappers}

We now prove the exclusion used in constructing the payload spaces.  Throughout
this section, $S$ is a primary strongly connected component containing an odd
closed walk and a profile triple $(a_0,b,k_0)$ with all three vertices in $S$.
We construct a small logical calculus inside the given PTS.  Its judgments will
be genuine PTS judgments in one finite open context.  The later contradiction
will therefore fall under the system-wide weak-normalization hypothesis.

We say that a type is \emph{sorted at a profile} $I$ if it has each sort in $I$.
Its actual profile may be larger than $I$.  All constructions below use only
these separate sort judgments, and never uniqueness of sorts or product rules.
Since $S$ is strongly connected and has an odd closed walk, any two of its
vertices are joined inside $S$ by paths of both parities: join the first vertex
to the walk and the walk to the second, and either insert the odd walk or omit
it.  A path's parity is the number of its negative edges modulo two.

\subsection{Modes and the ambient context}
\label{subsec:enc-modes}

Partition variable names into two infinite classes, of modes $\mathsf d$
and $\mathsf p$.  Sorts have mode $\mathsf d$.  Products, abstractions and
applications carry a label
\[
 (i,j)\in\{(\mathsf d,\mathsf d),(\mathsf d,\mathsf p),
                 (\mathsf p,\mathsf p)\}.
\]
In a binder, $i$ is the bound variable's mode.  A product
$\Pi^{i,j}x:A.B$ has mode $\mathsf d$ and both $A,B$ have mode
$\mathsf d$.  An abstraction $\lambda^{i,j}x:A.M$ has annotation of mode
$\mathsf d$ and body and result of mode $j$.  In an application with label
$(i,j)$, the argument has mode $i$ and the function and result have mode $j$.
Consequently, an expression of mode $\mathsf d$ has no free variable of mode
$\mathsf p$.  Declaration types and expected types always have mode
$\mathsf d$.

The typing rules are the original PTS rules with the displayed label retained
on the product in abstraction and application rules.  A beta step contracts
an abstraction and application with matching labels; reduction is closed under
all syntax positions, including annotations.  Conversion is generated by this
labelled reduction.  We suppress application labels only when the displayed
function type determines them.  The labelled metatheory and its relation to
ordinary reduction are proved in Lemma~\ref{lem:enc-labelled-metatheory}.
In particular, the weakening, substitution, conversion and generation
arguments used below apply in these modes.

\subsection{Reduction after forgetting the modes}
\label{sec:normal-erasure}

Before constructing the witnesses and wrappers, we establish the
labelled metatheory and its relation to the given pure type system.  Write $|M|$ for the term obtained by
forgetting every mode label.  Erasure keeps all applications, binders, and
annotations.  In particular, a redex in an annotation remains a redex after
erasure.

Labelling and erasure also occur in the structural transformations of
Roux and van Doorn \cite[Sections~3--4]{RouxVanDoorn2014}.
Here the exact interface is preservation of typing together with lifting
of every erased beta step. Keeping annotations is part of that interface;
the domain-erasure comparison studied by Barthe and Coquand
\cite[Theorem~17]{BartheCoquand2006} concerns a different operation.

\begin{lemma}[Labelled metatheory and reduction lifting]
\label{lem:enc-labelled-metatheory}
The labelled typing rules have context validity, correctness of types,
weakening, typed substitution by terms of the declared modes, context
conversion, generation, and subject reduction at the exact expected type.
Labelled beta reduction is confluent, and convertible
labelled products have the same pair of labels and convertible corresponding
components.  Erasing a labelled typing derivation gives a typing derivation
in the original pure type system.

If $\Gamma\vdash M:A$ is a labelled judgment and
$|M|\longrightarrow_\beta N$, there is a labelled $M'$ with
\[
 M\longrightarrow_\beta M',\qquad |M'|=N,
 \qquad \Gamma\vdash M':A,
\]
up to alpha conversion.  The mode of $M'$ is the mode of $M$.
Consequently, every finite reduction of $|M|$ lifts to a typed labelled
reduction.  If $|M|$ is weakly normalizing, some labelled reduct of $M$ has
beta-normal erasure.
\end{lemma}

\begin{proof}
Substitution preserves modes by induction on syntax.  A data expression has
no free proof variable, so substitution for a proof variable cannot introduce
proof syntax into a data expression.  Substitution for a data variable uses
a data term.  At binders, choose fresh names in the same mode class.

Context validity follows by rule induction, retaining the sort typing of
each declaration from its introduction premise.
The proofs of weakening and typed substitution are inductions on typing
derivations.  Each binder retains its label, and each substituted term has
the mode required by that binder or declaration.  Context conversion follows
by the identity substitution, typing the changed variable at its old type
by conversion.  Generation follows by removing final conversions and
weakenings.  None of these arguments compares distinct possible sort
assignments.

Correctness of types is another rule induction: an expected type is a literal
sort or is itself
sorted.  At application, the function's product type is sorted by induction;
generation on that product and typed substitution sort the application's
result type.  This remains true through conversion and weakening.

For confluence, use parallel reduction with homomorphic clauses at every
constructor, including annotation children, and a contracting clause only
when the application and lambda labels match.  Parallel substitution is
proved by syntax induction.  Every parallel reduct of a term parallel-reduces
to its complete development: contract the matching redexes already present
in the original term and develop all children.  At an unmatched
lambda--application pair only the homomorphic clause applies.  This proves
the triangle property, hence confluence.  The reflexive transitive closures
of parallel and ordinary reduction agree, since parallel contractions can be
performed from the children upwards.  Reduction never changes a product's
outer constructor or labels.  Common reducts therefore give labelled product
compatibility, and a product cannot be convertible to a sort.

Subject reduction now follows by induction on a typing derivation.  At a root
beta step, generation and product compatibility identify the lambda's
annotation with the application's domain up to conversion.  Convert the
argument to the annotation type, substitute in the body, and convert the
result to the original expected type.  A reduction of an application
argument changes its substituted result type by conversion.  A reduction
of a product domain is handled by context conversion in the codomain
premise.  For a lambda annotation step, reduce the corresponding domain in
its supporting product judgment, context-convert its body premise,
re-form the lambda, and convert back to the original sorted product type.
The other positions, including positions inside
annotations, follow by induction.  These constructions keep every label.
Erasure of typing follows directly by induction on the labelled derivation.

It remains to check that an erased redex can be contracted with the labels
present.  Every syntactic subterm of a typable term has a typing in the
appropriate binder context; for an annotation this follows from context
validity.  Consider the subterm at the erased redex.  Its shape is
\[
 (\lambda^{i,j}x:D.P)\mathbin{@^{i',j}}Q.
\]
The result modes already force the displayed common $j$.  Lambda generation
gives a product type with labels $(i,j)$; application generation gives a
convertible product type with labels $(i',j)$.  Product compatibility forces
$i=i'$.  Thus this is a matching labelled redex.  Contract it at the same
position.  Capture-free substitution commutes with erasure, so its erasure
is $N$; subject reduction gives the asserted judgment and substitution
preserves its mode.  This applies to any position, including one in a
lambda or product annotation.  Induction on a finite erased reduction proves
the remaining assertions.
\end{proof}

\subsection{Witnesses and terminal transformations}
For each feasible profile $I$ needed in the finite construction, choose a
normal witness $D_I$ sorted at $I$, using only data names and labels
$(\mathsf d,\mathsf d)$.  Such witnesses follow from the profile-generation
construction: literals realize nonempty axiom profiles, variables of sorted
literal type realize the singleton generators, and dummy products realize
$\operatorname{Out}(I,J)$.  Combine the finitely many witness contexts by
renaming their data variables disjointly and weakening.  Extend the result
by a data variable $e_I:D_I$ for each needed witness.  Denote this finite
valid data context by $\Delta_{\mathsf d}$.

Here and below every profile triple means
$K=\operatorname{Out}(I,J)$.  If $T$ is sorted at $J$, the positive edge
$J\longrightarrow K$ associated with this triple sends $T$ to
\[
   \Pi^{\mathsf d,j}z:D_I.T,
   \qquad j=\mathsf p\text{ for proof values},\quad
          j=\mathsf d\text{ for data values},
\]
where $z$ is fresh and absent from $T$.  This type is sorted at every sort
$c\in K$: choose $a\in I,b'\in J$ with $(a,b',c)\in\mathcal R$ and use
those two available sort judgments.  The corresponding lift and projection
are
\[
    \operatorname{up}(t)=\lambda^{\mathsf d,j}z:D_I.t,
    \qquad \operatorname{pr}(w)=w\,e_I.
\]
They have the displayed types and
$\operatorname{pr}(\operatorname{up}(t))=_\beta t$.
A positive inclusion edge $I'\longrightarrow I$ with $I\subseteq I'$
uses the identity type, lift and projection.

A negative edge $I\longrightarrow K$ associated with $(I,J,K)$ sends a
type $T$ sorted at $I$ to
$\Pi^{j,j}z:T.F$, where $F$ is a fixed type sorted at $J$ and $j$ is
$\mathsf p$ or $\mathsf d$ according to the construction.  Formation uses
exactly the same argument for each $c\in K$.  Thus a path $\rho$ defines
a wrapper $W_\rho(T)$ by successive applications of these operations,
with specified tails at negative edges.  Empty paths are permitted, and
\[
     W_{\rho\sigma}(T)=W_\sigma(W_\rho(T))
\]
by the definition.  We always specify whether the wrapper carries proof
values or data values.

For every internal negative edge of $S$ with triple $(I,J,K)$, fix the
proof tail $F_J=D_J$ before choosing any paths. Call these finitely many types the
\emph{terminals}; include $F_b=D_b$ by the negative edge of
$(a_0,b,k_0)$.  Each terminal is normal and has only
$(\mathsf d,\mathsf d)$ labels.

\begin{lemma}[Terminal transformations]
\label{lem:enc-terminal-converters}
There is a finite valid extension $\Delta=\Delta_{\mathsf d},
\Delta_{\mathsf p}$ and, for every pair of terminals $F,F'$, a proof
term $\tau_{F,F'}(p):F'$ in $\Delta,p:F$.  Every declaration of
$\Delta_{\mathsf p}$ has the form
\[
       c:\Pi^{\mathsf p,\mathsf p}u:T.F_J,
\]
where $T$ is a normal, constant telescope of
$(\mathsf d,\mathsf p)$ products ending at a terminal.  Its suffixes
are sorted, and its type contains no free proof variables.
\end{lemma}
\begin{proof}
Associate each terminal with a negative edge producing it.  Lift the input
terminal along the positive edge of its associated triple, reaching $S$.
Choose a path in $S$ to the domain vertex of an edge producing the target
terminal, and append that negative edge.  Carry a proof along this route.
At a positive edge use the lift above.  At a negative edge $(I,J,K)$,
if the currently carried type is $T$, introduce a fresh proof variable
$c:\Pi^{\mathsf p,\mathsf p}u:T.F_J$ and replace the carried proof $t$
by $c\,t:F_J$.  Unless this is the last edge, lift $c\,t$ along the
positive edge $J\longrightarrow K$ of this same triple before continuing.
The product type of $c$ is sorted at $K$, so the declaration is valid.

Initially the carried type is a terminal.  Positive edges add constant
$(\mathsf d,\mathsf p)$ prefixes, and negative edges reset the carried
type to a terminal before the prescribed lift.  Thus every converter domain
has precisely the asserted form.  Its normality and the sorting of each
suffix follow from its construction.  The types depend on chosen types and
paths, not on the input proof.  Choose the finitely many pairwise routes
first and predeclare all their converters in $\Delta_{\mathsf p}$;
weakening makes every resulting transformation available in the same
context.  Types use data only, so the converter declarations do not depend
on one another.
\end{proof}

\subsection{Proof transport and derived logic}
\label{subsec:enc-proof-transport}

Let $\rho$ be a proof path inside $S$ from a profile at which $T$ is sorted, and put
$W=W_\rho(T)$.  The following operations are defined in every valid
extension of $\Delta$.  All displayed inputs after a semicolon describe
an open term under the indicated fresh variable.
For an even path they have types
\[
  \operatorname{in}_\rho(p:T):W,
  \qquad \operatorname{out}_\rho(w:W;e.k):F
       \quad(e:T\vdash k:F),
\]
and for an odd path they have types
\[
  \operatorname{in}_\rho(e.k):W
       \quad(e:T\vdash k:F),
  \qquad \operatorname{out}_\rho(w:W,p:T):F,
\]
where $F$ is any terminal.  Path and terminal subscripts will often be
omitted, but remain part of each chosen operation.
For the empty path, $\operatorname{in}(p)=p$ and
$\operatorname{out}(w;e.k)=k[e:=w]$.
Appending a positive edge lifts the previous $\operatorname{in}$ and
projects the input to the previous $\operatorname{out}$.
For a negative edge with tail $F'$ appended to an even path $\rho$, set
\begin{align*}
 \operatorname{in}_{\rho-}(e.k)
    &=\lambda^{\mathsf p,\mathsf p}u:W_\rho(T).
        \operatorname{out}_\rho(u;e.\tau_{F,F'}(k)),\\
 \operatorname{out}_{\rho-}(w,p)
    &=\tau_{F',F}(w\,\operatorname{in}_\rho(p)).
\end{align*}
For a negative edge appended to an odd path $\rho$, set
\begin{align*}
 \operatorname{in}_{\rho-}(p)
    &=\lambda^{\mathsf p,\mathsf p}u:W_\rho(T).
        \operatorname{out}_\rho(u,p),
       &&\text{with output terminal }F',\\
 \operatorname{out}_{\rho-}(w;e.k)
    &=\tau_{F',F}(w\,\operatorname{in}_\rho(e.k)).
\end{align*}
These are recursion equations on path length.  Each abstraction uses the
already formed negative-edge product.  Each application uses its domain
exactly, and the terminal transformations have the displayed input and
output types.  This proves the typings by induction; callbacks are thinned
by insertion before introducing each fresh variable.

Call a type \emph{stable} if it is a syntactic telescope of proof-result
products, with labels $(\mathsf d,\mathsf p)$ or
$(\mathsf p,\mathsf p)$, ending at a terminal.  A telescope of length
zero is allowed.  Let $Q$ be a sorted stable type independent of a proof
variable $e:T$.  The even-path operation $\operatorname{out}(w;e.k)$
extends from terminal-valued callbacks to $k:Q$: introduce the successive
raw arguments $\vec x$ of $Q$, thin $w$ and the callback into that
context, apply terminal-valued $\operatorname{out}$ to
$e.k\,\vec x$, and abstract the same arguments $\vec x$.
Generation on the sorted telescope provides the domain sortings and the
supporting product judgments for all these abstractions.  In particular,
if $T$ is stable, the identity callback gives
\[
                    \operatorname{down}_\rho(w):T
                 \qquad (\rho\text{ even}).
\]
A proof wrapper preserves stability of its input.  A proof path containing
a negative edge produces a stable type even if its input was not stable,
because the first such edge ends at a terminal.

A \emph{formula} will mean a stable data type sorted at every sort of
$b$. A formula is therefore a data-mode expression used as a type; a proof
of that formula means a proof-mode inhabitant. The mode of an inhabitant
is not determined by its expected type.  The next operations give a typed classical logical calculus of
formulas.  These operations are defined terms, not additional PTS rules.

Fix even proof paths $\alpha:b\to a_0$, $\beta:b\to b$, and
$\gamma:k_0\to b$.  Define
\[
 P\Rightarrow Q=
 W_\gamma\bigl(\Pi^{\mathsf p,\mathsf p}p:W_\alpha(P).
                                      W_\beta(Q)\bigr).
\]
The product rule is supplied by $(a_0,b,k_0)$, and all wrappers are
already typed.  Given $q:Q$ in a context extended by $p:P$, define
\begin{align*}
 \lambda^*p.q
   &=\operatorname{in}_\gamma\bigl(
       \lambda^{\mathsf p,\mathsf p}p':W_\alpha(P).
       \operatorname{in}_\beta(q[p:=\operatorname{down}_\alpha(p')])
                                      \bigr),\\
 f\cdot p
   &=\operatorname{down}_\beta\bigl(
       \operatorname{down}_\gamma(f)\,
                          \operatorname{in}_\alpha(p)\bigr).
\end{align*}
They have types $P\Rightarrow Q$ and $Q$, respectively.  Data formulas
cannot contain $p$, so the product codomain is independent of this proof
variable.

For a data type $T$ sorted at a direct vertex $I$, fix once and for all
a triple $(I,J,K)$ with $J,K\in S$, and even proof paths
$\alpha_I:b\to J$ and $\gamma_I:K\to b$.  For a formula $Q$ in
context $x:T$, put
\[
  \forall x:T.Q=
  W_{\gamma_I}\bigl(\Pi^{\mathsf d,\mathsf p}x:T.
                                       W_{\alpha_I}(Q)\bigr).
\]
The associated introduction and elimination are
\[
 \Lambda x.q=\operatorname{in}_{\gamma_I}
       (\lambda^{\mathsf d,\mathsf p}x:T.
                                  \operatorname{in}_{\alpha_I}(q)),
 \qquad
 f[t]=\operatorname{down}_{\alpha_I}
                   (\operatorname{down}_{\gamma_I}(f)\,t).
\]
They are typed using precisely the chosen triple.  The direct vertex and
paths are part of the domain annotation and remain fixed under
substitution; paths are never reselected from a newly enlarged profile.
These constructions show that implication and universal quantification
again yield formulas.

Set $\bot=F_b$, $\neg P=P\Rightarrow\bot$, and
$\top=\bot\Rightarrow\bot$.  There are derived operations
\[
        \operatorname{abs}_P:\bot\longrightarrow P,
        \qquad \operatorname{dn}_P:\neg\neg P\longrightarrow P,
\]
where the arrows describe input and output judgments, not an assertion that
these are primitive function types.  To construct either operation, expose
the raw telescope arguments $\vec x$ of $P$, ending at a terminal $F$.
For $f:\bot$, abstract the terminal body $\tau_{\bot,F}(f)$.
For $w:\neg\neg P$, abstract the terminal body
\[
 \tau_{\bot,F}\left(
       w\cdot(\lambda^*p.\tau_{F,\bot}(p\,\vec x))\right).
\]
The applications $p\,\vec x$ use the original labels and annotations
of $P$.  Sorting and abstraction support follow from generation on $P$,
as in the definition of $\operatorname{down}$.

For completeness, the other logical operations used below are the
following explicit abbreviations and term constructors:
\begin{align*}
 P\land Q&=\neg(P\Rightarrow\neg Q),&
 (p,q)_\land&=\lambda^*h.(h\cdot p)\cdot q,\\
 \operatorname{fst}(w)
   &=\operatorname{dn}_P
       (\lambda^*n.w\cdot(\lambda^*p.\lambda^*q.n\cdot p)),&
 \operatorname{snd}(w)
   &=\operatorname{dn}_Q
       (\lambda^*n.w\cdot(\lambda^*p.\lambda^*q.n\cdot q)),\\
 P\leftrightarrow Q&=(P\Rightarrow Q)\land(Q\Rightarrow P),&
 \operatorname{to}(e,p)&=\operatorname{fst}(e)\cdot p,\\
 &&\operatorname{back}(e,q)&=\operatorname{snd}(e)\cdot q,\\
 P\lor Q&=\neg(\neg P\land\neg Q),&
 \operatorname{inl}(p)&=\lambda^*h.\operatorname{fst}(h)\cdot p,\\
 &&\operatorname{inr}(q)&=\lambda^*h.\operatorname{snd}(h)\cdot q.
\end{align*}
Given $w:P\lor Q$ and branches $p:P\vdash u:R$ and
$q:Q\vdash v:R$, their case term is
\[
 \operatorname{dn}_R\left(
   \lambda^*n.w\cdot
       (\lambda^*p.n\cdot u,\lambda^*q.n\cdot v)_\land\right).
\]
Excluded middle has the term
\[
      \operatorname{em}(P)=
        \lambda^*h.\operatorname{snd}(h)\cdot\operatorname{fst}(h)
          :P\lor\neg P.
\]
For the same permitted quantified domains as above, put
\[
 \exists x:T.Q=\neg\forall x:T.\neg Q.
\]
If $t:T$ and $p:Q[x:=t]$, a witness term is
$\operatorname{wit}(t,p)=\lambda^*h.h[t]\cdot p$.
If $w:\exists x:T.Q$ and $u:R$ under $x:T,p:Q$, where $R$ is
independent of $x,p$, existential elimination is
\[
 \operatorname{open}_R(w;x,p.u)=
  \operatorname{dn}_R\bigl(
    \lambda^*n.w\cdot(\Lambda x.\lambda^*p.n\cdot u)\bigr).
\]
Each display can be checked from the introduction and elimination
judgments already proved.  All new assumptions are introduced with their
formula or data type as annotation, and terms are weakened when necessary.

\begin{lemma}[Finite propositional derivations]
\label{lem:enc-propositional}
If formulas $P_1,\ldots,P_m$ propositionally entail a formula $R$, there
is a specified finite term
$\mathsf{pl}_R(p_1,\ldots,p_m):R$ under $p_i:P_i$.
Here $\bot$ has truth value false; specified Boolean connectives are
parsed, and other subformulas, including quantifiers, may be treated as
atoms.
\end{lemma}
\begin{proof}
List the finitely many atoms and nest cases on their terms
$\operatorname{em}$.  Each resulting branch supplies either the atom or
its negation.  Recursively build a proof of every true parsed formula and
a proof of the negation of every false parsed formula.  For $\bot$ the
required negation is the identity term.  For conjunction, a true case
uses pairing, and a false case contradicts the false conjunct's projection.
For disjunction, a true case uses its true injection, and a false case
uses case analysis to contradict either disjunct.  For implication, a true
case either returns a proof of the consequent or derives it by
$\operatorname{abs}$ from the false antecedent; a false case contradicts
the application of an assumed implication to the already proved antecedent.
Negation is implication to $\bot$, and biconditional is its displayed
conjunction.  If some premise is false, apply its constructed negation to
the given proof and use $\operatorname{abs}_R$.  Otherwise the assumed
truth-table entailment makes $R$ true, so use its constructed proof.
The recursively typed case terms discharge all branch assumptions.
\end{proof}

We also use three explicit congruence constructors.  For a proof
$e:P\leftrightarrow Q$ under data $x:T$, define
\[
 \mathsf{All}_x(e)=
 \bigl(\lambda^*h.\Lambda x.\operatorname{to}(e,h[x]),
       \lambda^*h.\Lambda x.\operatorname{back}(e,h[x])\bigr)_\land.
\]
It compares $\forall x:T.P$ and $\forall x:T.Q$ with exactly the same
annotated domain.  Define $\mathsf{Ex}_x(e)$ by pairing the two functions
which open an existential input, retain its data witness, and use
$\operatorname{to}$ or $\operatorname{back}$ on its proof witness.
For a proof $e:P\leftrightarrow Q$ under $v:D$, where $D$ is a formula,
put
\[
 \mathsf{Gd}_{v:D}(e)=
 \bigl(\lambda^*h.\lambda^*v.\operatorname{to}(e,h\cdot v),
       \lambda^*h.\lambda^*v.\operatorname{back}(e,h\cdot v)\bigr)_\land.
\]
This compares $D\Rightarrow P$ and $D\Rightarrow Q$.
Multiple indices mean iteration of the corresponding constructor.

These conventions make later derivations finite algorithms for labelled
terms.  A stated propositional consequence invokes
Lemma~\ref{lem:enc-propositional}; proving the two directions of an
 equivalence means abstracting the two assumptions and pairing; opening an
existential means the displayed $\operatorname{open}$ term at the stated
target.  A schematic data parameter is an open parameter, not an implicit
universal quantifier.  Only displayed $\Lambda$ operations quantify over
data.  Logical equivalences are used to construct proofs and never to
replace a data subterm.  All fixed paths, roles and syntax choices are
preserved under substitution.

\subsection{Data wrappers and exact cancellation}
\label{subsec:enc-data-transport}

For a sort $u$ with $\{u\}\in S$ and
$\operatorname{Ax}(u)\ne\varnothing$, fix an even proof path from $b$
to $\{u\}$.  For a formula $P$, its wrapper type along this path is a
data expression of type $u$; denote it by $\operatorname{encode}_u(P)$.
Fix also an even proof path from $\{u\}$ to $b$ containing a negative
edge.  Such a path is obtained from any even one by inserting an odd
closed detour twice.  Wrapping a data expression $t:u$ along this second
path defines a stable formula $\operatorname{decode}_u(t)$.
The two proof-path inclusions and the two stable retractions give a proof
\begin{equation}
 \operatorname{decode}_u(\operatorname{encode}_u(P))
                             \leftrightarrow P.
 \label{eq:enc-code-decode}
\end{equation}
Explicitly, the forward implication applies the two $\operatorname{down}$
operations and the reverse implication applies the two
$\operatorname{in}$ operations.  No equality of the two formulas is asserted.

Recall that $H$ is the positive-path closure of the nonempty
$\operatorname{Ax}(u)$ with $\{u\}\in S$.  For each tail vertex $J$
of a retained negative secondary edge, choose such a $u$ and an all-positive
primary path $\eta:\operatorname{Ax}(u)\to J$.  Define the data type
\[
                 U_J=W_\eta(u),
\]
using data labels.  It is sorted at $J$.  Define
$\operatorname{send}_J(P):U_J$ by positively lifting
$\operatorname{encode}_u(P):u$, and define
$\operatorname{read}_J(t)$ by positively projecting $t:U_J$ to type
$u$ and applying $\operatorname{decode}_u$.  Positive
projection--lift beta conversion and \eqref{eq:enc-code-decode} give a
specified proof
\begin{equation}
        \operatorname{read}_J(\operatorname{send}_J(P))
                                     \leftrightarrow P.
        \label{eq:enc-read-send}
\end{equation}
These choices are fixed once for each retained edge.

Now use data wrappers on secondary paths, and on all-positive primary
paths when indicated.  At a retained negative edge $(I,J,K)$, the tail is
$U_J$ and the label is $(\mathsf d,\mathsf d)$.  For an even data path
$\rho$ from the sorting profile of $T$, define
\[
 \operatorname{pack}_\rho(t:T):W_\rho(T),\qquad
 \operatorname{use}_\rho(w:W_\rho(T);x.\Phi)\quad\text{a formula},
\]
where $\Phi$ is a formula under data $x:T$.  For an odd data path define
\[
 \operatorname{mk}_\rho(x.\Phi):W_\rho(T),\qquad
 \langle w:W_\rho(T),t:T\rangle_\rho\quad\text{a formula}.
\]
At the empty path use $\operatorname{pack}(t)=t$ and
$\operatorname{use}(w;x.\Phi)=\Phi[x:=w]$.
A positive extension uses the same lift and projection operations as
before.  If a negative edge with tail $U_J$ is appended to an even path,
put
\begin{align*}
 \operatorname{mk}_{\rho-}(x.\Phi)
   &=\lambda^{\mathsf d,\mathsf d}v:W_\rho(T).
           \operatorname{send}_J(\operatorname{use}_\rho(v;x.\Phi)),\\
 \langle w,t\rangle_{\rho-}
   &=\operatorname{read}_J(w\,\operatorname{pack}_\rho(t)).
\end{align*}
If it is appended to an odd path, put
\begin{align*}
 \operatorname{pack}_{\rho-}(t)
   &=\lambda^{\mathsf d,\mathsf d}v:W_\rho(T).
                         \operatorname{send}_J(\langle v,t\rangle_\rho),\\
 \operatorname{use}_{\rho-}(w;x.\Phi)
   &=\operatorname{read}_J(w\,\operatorname{mk}_\rho(x.\Phi)).
\end{align*}
All expressions in these definitions have mode $\mathsf d$.
Path induction proves their typings: the body of each new lambda has type
$U_J$, the negative-edge product is sorted by $(I,J,K)$, and each
application supplies an argument at exactly its displayed domain.  The
formula returned by a read is stable by the negative-containing decode
path.  Fresh variables and weakening keep callbacks in scope.

\begin{lemma}[Exact data cancellation]
\label{lem:enc-data-cancellation}
For an even data path there is a proof of
\[
       \operatorname{use}_\rho(\operatorname{pack}_\rho(t);x.\Phi)
                             \leftrightarrow\Phi[x:=t],
\]
and for an odd data path there is a proof of
\[
             \langle\operatorname{mk}_\rho(x.\Phi),t\rangle_\rho
                             \leftrightarrow\Phi[x:=t].
\]
These are uniform term constructors in all the displayed open data
parameters.  The right-hand sides contain the original input $t$.
\end{lemma}
\begin{proof}
Induct on the path.  The empty-path equivalence is propositional
reflexivity.  A positive extension beta-reduces its projection of a lift,
so conversion reduces its assertion to the preceding one.  For a negative
extension, the displayed lambda beta-reduces the left side to
$\operatorname{read}_J(\operatorname{send}_J(\Psi))$, where $\Psi$
is the left side of the previous cancellation assertion.  Apply
\eqref{eq:enc-read-send}, the induction hypothesis, and
$\mathsf{pl}$ for transitivity of equivalence.  The expansions commute
with capture-free substitution because all paths, tails and names are
fixed hygienically.  Consequently the final callback is exactly
$\Phi[x:=t]$ up to labelled conversion, with no extensional substitution
inside data.
\end{proof}

The preceding constructions use finitely many sorts, profiles, witnesses,
paths and converters.  Each term in the remainder is built from finitely
many of their instances.  Thus all the resulting judgments have the one
finite ambient context $\Delta$, with any explicitly stated local
parameters.  The next construction uses these operations to encode the
forbidden profile configuration.
